Encoding in SAT can be \textbf{challenging}. If we are not careful about the encoding phase we may encounter big problems: we cannot directly \textbf{translate}
any \textbf{Constraint Programming problem} into a \textbf{Satisfiability problem}, as we already did for N-Queens problem in the previous chapter. \vspace{7pt}

The cardinality of the boolean decision variables used, as the number of clauses, can depend also from the dimension of the problem. An approach in which we have 
a total amount of variables and clauses similar to $n^2$ or worse is not a suitable approach for SAT problems. Not only the resolution of the problem 
might be very slow, but also the translation of the original propositional formula to its CNF form can be heavy.